% SPDX-License-Identifier: Apache-2.0
%
% Formal verification in TxHDL: the tutorial on proving what a unit
% states about itself. Built by //docs:prove. Every listing is included
% from the tree and every printout, log and figure is what the build
% made by running the example and the proofs.
\documentclass[journal,9pt]{IEEEtran}

\newcommand{\listingsize}{\fontsize{6}{7}\selectfont}
% SPDX-License-Identifier: Apache-2.0
%
% The house preamble, shared by every document in this directory. Loaded
% after \documentclass, before \title. Keeping it in one file is what
% stops two articles drifting apart in font, listing style or figure
% style.
% Computer Modern. IEEEtran selects Times (ptm) by default, so the face has
% to be chosen explicitly.
%
% Latin Modern is the Computer Modern variety used here, and the choice is
% forced rather than aesthetic. Under T1 encoding, plain `cmr` has no Type 1
% outlines unless cm-super is installed, and it is not in the pinned
% distribution, so pdflatex falls back to EC bitmaps and every glyph in the
% PDF becomes a Type 3 font. That was measured: the first build produced a
% document with 10 Type 3 fonts and no Type 1. Latin Modern is Computer
% Modern redrawn with a full T1 repertoire and real .pfb outlines, so the
% same design comes out scalable.
\usepackage[T1]{fontenc}
\usepackage[utf8]{inputenc}
\usepackage{lmodern}
% microtype is not in the pinned TeX distribution. emergencystretch alone
% keeps the two-column measure from producing overfull lines.
\setlength{\emergencystretch}{3em}

\usepackage{amsmath}
\usepackage{amssymb}
\usepackage{amsthm}
\usepackage{booktabs}
\usepackage{array}
% enumitem is not in the pinned TeX distribution; the list spacing below
% does what \setlist would have done.
\usepackage{float}
\usepackage{xcolor}
\usepackage{listings}
\usepackage{tikz}
\usetikzlibrary{arrows.meta,positioning,fit,backgrounds,calc}
\usepackage[hidelinks,bookmarks=true]{hyperref}

% Numbered, definition-style box for a rule the reader is meant to apply.
\theoremstyle{definition}
\newtheorem{rulebox}{Rule}
\newtheorem{example}{Example}

% Unnumbered asides.
\theoremstyle{remark}
\newtheorem*{tip}{Tip}
\newtheorem*{pitfall}{Pitfall}

% Tighter lists, in place of enumitem's \setlist.
\setlength{\topsep}{2pt}
\setlength{\itemsep}{1pt}
\setlength{\parsep}{0pt}

\newcommand{\code}[1]{\texttt{#1}}
\newcommand{\term}[1]{\emph{#1}}

% Syntax highlighting. Four roles, distinguishable in colour and still
% distinguishable in greyscale, because an IEEE article gets printed: the
% keyword is the only bold role, the comment is the only italic one, and the
% two remaining roles differ in lightness as well as in hue.
\definecolor{kwblue}{RGB}{20,60,150}
\definecolor{cmtgreen}{RGB}{90,120,90}
\definecolor{strred}{RGB}{150,50,40}
\definecolor{numgray}{gray}{0.45}
\definecolor{framegray}{gray}{0.70}
\definecolor{bgshade}{gray}{0.97}
% A darker fill for figure cells. The listing background has to stay very
% light to keep code readable; a figure cell has to be clearly shaded or the
% distinction it draws is invisible in print.
\definecolor{cellshade}{gray}{0.90}

% The listing size is the document's choice, because it depends on the
% column: a two-column article sets `\listingsize` to 6pt and keeps its
% listings within 70 columns, a one-column document keeps the body size
% and includes source files at 80. Lines are never broken; a line that does not fit
% overflows the frame, which is visible, rather than wrapping, which is
% not.
\providecommand{\listingsize}{\scriptsize}

% `'` is deliberately not a string delimiter: the language uses it for loop
% labels, as Rust does, so treating it as a quote would swallow the rest of
% the line.
\lstdefinestyle{txhdl}{
  basicstyle=\ttfamily\listingsize,
  keywordstyle=\bfseries\color{kwblue},
  commentstyle=\itshape\color{cmtgreen},
  stringstyle=\color{strred},
  numberstyle=\tiny\color{numgray},
  morekeywords={mod,use,pub,const,type,struct,enum,trait,impl,fn,async,let,mut,
    match,if,else,loop,while,for,in,break,continue,return,as,await,
    module,bus,channel,transaction,protocol,pipeline,flow,seq,config,
    port,stage,pipe,instance,connect,bind,tag,domain,blackbox,foreign,
    property,role,signal,handler,true,false,
    sequence,interface,modport,implements,package,import,var,wire,func,proc,
    case,default,next,fork,branch,join,state,fsm},
  morecomment=[l]{//},
  morestring=[b]",
  showstringspaces=false,
  columns=fullflexible,
  keepspaces=true,
  breaklines=false,
  % Framed, so a listing is bounded rather than floating in the column.
  frame=single,
  rulecolor=\color{framegray},
  framesep=3pt,
  backgroundcolor=\color{bgshade},
  xleftmargin=3pt,
  xrightmargin=3pt,
  aboveskip=6pt,
  belowskip=6pt,
  % Every listing is numbered and captioned, so it can be referred to by
  % number. `caption` without `float` numbers the listing and leaves it
  % where it was written, which is what a listing in running prose needs.
  captionpos=b,
  abovecaptionskip=2pt,
  belowcaptionskip=6pt,
}

% Figures drawn as text. Framed the same way, and with the highlighting
% switched off, because a box-and-arrow diagram is not code and half its
% words turning blue reads as an accident. Setting a style adds keys rather
% than clearing the ones already set, so the three roles are neutralised by
% hand rather than by leaving them out.
\lstdefinestyle{ascii}{
  basicstyle=\ttfamily\listingsize,
  keywordstyle=\ttfamily,
  commentstyle=\ttfamily,
  stringstyle=\ttfamily,
  columns=fullflexible,
  keepspaces=true,
  frame=single,
  rulecolor=\color{framegray},
  framesep=3pt,
  backgroundcolor=\color{bgshade},
  xleftmargin=3pt,
  xrightmargin=3pt,
  aboveskip=6pt,
  belowskip=6pt,
}
\lstset{style=txhdl}

% House TikZ styles. `shorten` lives on the arrow styles rather than on each
% arrow, so every arrow in the document gets the gap, including ones added
% later. TikZ clips a path at a node's border, so without it an arrowhead
% butts against the box with no gap at all. Keep `node distance` well above
% twice the shorten, or the arrow shrinks to nothing.
\tikzset{
  box/.style={draw,rounded corners=2pt,align=center,inner sep=4pt,
              font=\footnotesize},
  boxf/.style={box,fill=bgshade},
  lbl/.style={font=\scriptsize\itshape,align=center},
  fl/.style={-{Stealth[length=4pt]},shorten >=2.5pt,shorten <=2.5pt},
  fld/.style={fl,dashed},
}



\InputIfFileExists{buildstamp.tex}{}{}
\providecommand{\buildstamp}{Draft build (outside the Bazel flow).}

\title{Formal Verification in TxHDL}

\author{Filip Filmar and Dragiša Janković%
\thanks{Both authors are independent. Filip Filmar is the
corresponding author, at f@hdlfactory.com; Dragiša Janković is
reached at linkedin.com/in/dragisa-jankovic.}%
\thanks{This is the tutorial on proving what a unit states about
itself, for a reader who has met TxHDL through the step-by-step
tutorial~\cite{tutorial} and wants a check to hold for every input
rather than for one run. Every listing is included from the project
repository \code{HDL/txhdl} at build time, and every printout, log
and figure is what the build produced by running the example and the
proofs. The example is in the examples document~\cite{examples}, and
the cover~\cite{cover} lists every document.}%
\thanks{The authors used a large language model, Claude, as an
assistant in exploring the concepts and in writing and constructing
the documents and the programs; every commit of the repository says
so and carries the prompts that led to it, verbatim.}%
\thanks{\buildstamp}}

\markboth{Formal Verification in TxHDL}{Filmar and Janković: Formal Verification in TxHDL}

\begin{document}
\maketitle

\begin{abstract}
A TxHDL unit can state what must hold of it, what it takes for
granted of its inputs, and what it should be able to reach, with
\code{check!}, \code{assume!} and \code{cover!}, where it is written.
The run checks those statements on the one run it makes, and the two
co-simulations check the netlist against that run. This document
takes the same statements to a proof: SymbiYosys, on the Yosys the
build already pins, proves the checks of a decimal digit for every
input and every reachable state, with no solver beyond the one inside
yosys-abc; a bounded search with z3 finds the cycle in which a
broken check fails and draws it; a cover run finds the path to the
cover point. The proof is a test, \code{bazel test //...} runs it in
about a second, and a broken copy of the unit fails it, which is what
makes the passing one worth anything.
\end{abstract}

\section{Three Ways of Being Right}
\label{sec:p-three}

Every example in the tree is checked by comparison. The build runs the
example, records every signal of the run in a trace, lowers the unit
to VHDL and to Verilog, and replays the trace against both under nvc
and under Verilator: the netlist must do what the run did, at every
tick. That catches a lowering that is wrong. It says nothing about a
design whose Rust is wrong in the same way as its netlist, and it says
nothing about an input the run did not give.

The second way is to state what must hold where it is written.
\code{check!(c, "why")} says that \code{c} holds on every edge at
which the statement is reached; \code{assume!(c, "why")} says what
the unit takes for granted of its inputs; \code{cover!(c, "why")}
says that \code{c} can happen. The run stops on a broken check, or on
an assumption its own inputs break, and counts each cover point. The
netlist carries the same statements as assertions of the clocked
process they are in, so both co-simulations check them too. The
tutorial~\cite{tutorial} and the examples document~\cite{examples}
say this much.

The third way is the subject of this document. A statement in the
netlist is a statement about a finite machine, and a formal tool can
decide it for every input and every state the machine can reach, not
for the inputs one run happened to give. What follows is one unit,
its statements, its run, and then its proof, a failure, and a cover,
each with what the tool wrote.

\section{The Unit}
\label{sec:p-unit}

\code{ex\_formal} is a decimal digit that counts by the step it is
given: nought, one or two, wrapping past nine. It states four things.
The count is always a digit. Under the condition that it wraps, it
lands on nought or one. Its step is never three, which it does not
handle: an assumption, since the step is an input. And the count
reaches nine: a cover point.

\lstinputlisting[rangebeginprefix=//\ begin\{,rangebeginsuffix=\},rangeendprefix=//\ end\{,rangeendsuffix=\},includerangemarker=false,linerange=unit-unit,caption={From \code{lib/examples/ex\_formal.rs}: the digit and its four statements.},label={lst:p-unit}]{lib/examples/ex_formal.rs}

The run drives the step through nought, one and two, sixteen times,
and prints the count after each edge, then how often nine was
reached, then the Verilog the lowering wrote. Every statement held on
every edge, and the netlist holds the statements as SystemVerilog
immediate assertions between \code{`ifdef FORMAL} and \code{`endif},
inside the reset's \code{else}, under the same conditions the source
put them under.

\lstinputlisting[style=ascii,caption={What \code{ex\_formal} printed when the build ran it: the run, the cover count, and the netlist the proof reads.},label={lst:p-out}]{out_formal.txt}

\begin{figure*}[t]
\centering
\input{formal_timing.tex}
\caption{The run: the digit counting by its step and wrapping past
nine. One run, sixteen edges, and the statements held on each of
them.}
\label{fig:p-run}
\end{figure*}

\section{The Proof}
\label{sec:p-proof}

The proof reads the Verilog of Listing~\ref{lst:p-out} with
\code{FORMAL} defined, so the assertions are in. SymbiYosys drives
Yosys through the preparation: \code{read -formal} reads the file with
the assertions as formal cells, \code{prep} elaborates the module, and
SymbiYosys's own passes turn the clocked immediate assertions into
properties of the state machine, which is the step Yosys's plain
\code{sat} command gets wrong on its own. Then an engine runs.

For the unbounded proof the engine is \code{abc pdr}: property
directed reachability, inside yosys-abc, which the build already
pins for the ASIC flow~\cite{cover}. It needs no SMT solver. It finds
an inductive invariant, a set of states that contains the initial
state, is closed under every step the assumptions allow, and lies
inside every check; when it finds one, the checks hold for ever. The
digit's invariant is three frames deep.

\lstinputlisting[style=ascii,caption={The proof's log, as the build wrote it: both checks proved in frame three, and no trace, since nothing failed.},label={lst:p-proof}]{formal_proof.log}

The two outputs the engine names are the two checks; the assumption
is the constraint the warning in the full log speaks of, and the
cover point is not part of a proof.

\section{A Broken Unit}
\label{sec:p-broken}

A proof that has never failed proves nothing about the tool. So the
build makes two broken copies of the netlist and proves those too,
expecting each to fail. One removes the assumption, so a step of
three is allowed and the count leaves the digits. The other tightens
the first check from \code{count <= 4'h9} to \code{count <= 4'h8},
which nine breaks. Here is the tightened one under the same engine.

\lstinputlisting[style=ascii,caption={The tightened check under \code{abc pdr}: a failure, with no trace, since the engine is asked for a verdict and not a path.},label={lst:p-tight}]{formal_tight_proof.log}

A verdict is what a test wants; a reader wants to see the cycle. For
that the bounded search runs: \code{smtbmc} with z3, twenty steps
deep, which asks at each step whether some input takes the machine
from its initial state to a state that breaks a check, and stops at
the first step where the answer is yes, with the inputs that do it.

\lstinputlisting[style=ascii,caption={The bounded search on the tightened check: assumptions and assertions checked step by step, the failure in step six, and the trace written.},label={lst:p-search}]{formal_tight_search.log}

The trace is a VCD, and the build draws it with the same tools that
draw a run. Figure~\ref{fig:p-cex} is the counterexample: steps of
two, two, one, two and two take the count from nought to nine in
five edges, and at the sixth the tightened check reads nine and
fails. No run of the example gave those steps; the solver did.

\begin{figure}[htbp]
\centering
\input{formal_cex_timing.tex}
\caption{The counterexample the bounded search found: the steps that
take the count to nine, where a check of \code{count <= 8} fails.}
\label{fig:p-cex}
\end{figure}

\section{The Cover Point}
\label{sec:p-cover}

A cover run asks the opposite question: not whether a state can be
reached that breaks a check, but whether one can be reached that
satisfies the cover point. The same bounded engine answers it, and
the answer is a path.

\lstinputlisting[style=ascii,caption={The cover run: the count reaching nine in step six, and the trace to it.},label={lst:p-cover}]{formal_covered.log}

\begin{figure}[htbp]
\centering
\input{formal_reach_timing.tex}
\caption{The path to the cover point: four steps of two, one of one,
and the count is nine.}
\label{fig:p-reach}
\end{figure}

\section{The Rule and the Tests}
\label{sec:p-rule}

All of this is one script and two Bazel macros in
\code{tools/formal}. \code{formal\_test} runs SymbiYosys on the
Verilog a \code{waveform()} wrote, in one of three modes, and passes
only when SymbiYosys exits the way the test expects: nought for a
pass, two for a failure. Nothing else passes, and in particular an
UNKNOWN, exit code four, never does; induction at a small depth
answers UNKNOWN for a real bug, and a rule that took only a failure
as a failure would pass broken code. \code{formal\_run} is the same
run as a build step, whose log and trace this document includes, and
\code{formal\_trace} draws the trace.

\lstinputlisting[rangebeginprefix=\#\ begin\{,rangebeginsuffix=\},rangeendprefix=\#\ end\{,rangeendsuffix=\},includerangemarker=false,linerange=formal-formal,caption={From \code{docs/BUILD.bazel}: the two broken copies, the seven tests, the four runs and the two drawings.},label={lst:p-build}]{BUILD.bazel}

The script is the whole of the mechanism, and it is short enough to
read. It finds the two Debian trees the build unpacked, Yosys and z3,
by a file each holds; it wraps Yosys's two Python helpers so that they
run on the hermetic Python with the tree's modules on their path,
since the scripts Debian ships name \code{/usr/bin/python3} and
\code{/usr/share/yosys}, neither of which is where the tree is; it
writes the \code{.sby} file for the mode asked for; and it compares
the exit code with the one expected.

\lstinputlisting[style=ascii,caption={\code{tools/formal/formal.sh}, which every proof in the tree runs through.},label={lst:p-script}]{tools/formal/formal.sh}

Nothing is installed. SymbiYosys is fetched from its release tag by
checksum and installed the way its own Makefile installs it, into
\code{//third\_party/sby}; \code{click}, the one package it imports
beyond the standard library, comes from a hashed lock on the
Python~3.12 toolchain the build already has; z3 is a Debian closure
under \code{//third\_party/z3}, pinned the way Yosys is; and the
proof itself needs none of z3, since PDR runs inside yosys-abc. The
seven tests take about a second each, so they are ordinary tests of
\code{bazel test //...} rather than manual ones, and this document is
built and released with the rest.

\section{What Is Not Here}
\label{sec:p-not}

A proof is of the netlist, under the assumptions the unit states. An
assumption that is too strong proves a check nobody wanted; the
assumption-removed copy is in the tests for the opposite reason, to
show that the assumption is doing work.

The netlist puts a unit's statements under the reset's \code{else}
only when it added the reset port itself. A unit that declares its
own \code{rst}, as the Vreteno core does, has its checks evaluated
during reset as well, and its first proof would fail there for no
real reason; that is issue~633, to be fixed before the core states
anything.

The next units to state things are the AXI-Lite peripherals, against
a protocol checker, and the core itself, against the RISC-V formal
interface; both are proofs of the same shape as this one, on a
netlist a few thousand times larger, and neither is in the tree yet.

\begin{thebibliography}{9}

\bibitem{tutorial}
F.~Filmar and D.~Janković, ``TxHDL, step by step,'' project
repository \code{HDL/txhdl}, \code{//docs:tutorial}, 2026.

\bibitem{examples}
F.~Filmar and D.~Janković, ``TxHDL by example,'' project repository
\code{HDL/txhdl}, \code{//docs:examples}, 2026.

\bibitem{cover}
F.~Filmar and D.~Janković, ``TxHDL documents,'' project repository
\code{HDL/txhdl}, \code{//docs:cover}, 2026.

\bibitem{sby}
YosysHQ, ``SymbiYosys (sby), front-end for Yosys-based formal
verification flows,'' \code{github.com/YosysHQ/sby}, release
\code{v0.52}, 2025.

\end{thebibliography}

\end{document}
